\subsection{Design contract: delay decisions, retain legality}
\label{sec:ir-design}
The layers answer different questions about the same program. Domain asks what
messages and reductions mean. Iter asks which relation coordinates must be
visited. Kernel asks how that work is mapped to execution resources. Task and
Storage ask when work may run and which physical copy supplies its operands.
Tensor IR is a numerical-expression layer used alongside these decisions, not
necessarily a final fifth stage. The split avoids encoding a GPU schedule in
the mathematical interface, but does not make scheduling automatic or optimal.

The key invariant is preservation of the observable outputs and supported VJPs
for the declared input contract. Shape, dtype, device, field role, reducer state,
topology snapshot and data-access effects are parts of that contract. A pass
that cannot establish a required property must reject the transformation or
select an explicitly diagnosed alternative. A topology version in the IR is a
validity annotation; it does not make arbitrary external in-place mutation safe.

\subsection{Domain: field roles and reducer algebra}
In the example, the source field is indexed by a source ID and the edge weight
by an edge position. Both are floating-point vectors, but exchanging their roles
changes the program. \code{gf.apply} therefore retains \code{input\_roles},
\code{input\_names}, reducer symbols, region kinds and input segments. Its region
contains the edge-local arithmetic, not an already materialized global message
array. A relation separately retains cardinalities, origin, lifecycle, version
and the available degree constraints.

A reducer describes identity, lift, combine and finalize regions, including the
types of its intermediate state. This is richer than attaching the string
``sum'' to an arbitrary scatter. A product or online-softmax reducer needs
different state and different differentiation rules. Empty rows start from the
identity and still obey finalize. Associativity admits regrouping; commutativity
admits reordering. These declarations are semantic obligations, not properties
the compiler proves for unrestricted user-written Python. Floating-point sums
require tolerance-aware validation when their reduction tree changes.

The full fixture and four unedited compiler outputs are included in
\code{generated/ir/}. This excerpt shows the fixture's edge region; its
enclosing Domain operation is omitted:
\begin{lstlisting}[language=TigaIR]
^bb0(%source: f32, %edge: f32):
  %message = arith.mulf %source, %edge : f32
  "gf.yield"(%message) : (f32) -> ()
\end{lstlisting}
The block arguments represent one logical edge invocation. The region itself
does not say which CPU core, CUDA thread or memory tier executes it.

This representation makes the materialization decision explicit. An eager
implementation gathers a source row for every edge, writes its weighted
message and then reduces those messages. A fused row traversal reads the same
source and weight, forms a tile-local message and immediately updates reducer
state. For feature width $F$, the former can allocate an edge-message array
with $EF$ entries; the latter needs only tile state in addition to graph,
inputs and the $NF$ output. The edge set and message arithmetic are unchanged.
The fixed-topology study in Section~\ref{sec:supplement} holds that edge set
constant to examine execution-path differences. Generated radius traversal
additionally avoids accepted-edge storage, but that distinct benefit is not
isolated by a stored-CSR comparison.

\subsection{Iter: coordinates before threads}
The Domain-to-Iter pass replaces the application with \code{gf\_iter.traverse}
and preserves the reducer and edge region. In the checked CSR fixture, the
attributes \code{coordinate\_hierarchy = "compressed-row"} and
\code{ordering = "destination-major"} describe a row-oriented coordinate walk.
These names do not mean that the operation is already an explicit sequential
loop, or that all threads use the same row length.

For the three-row example, \code{row\_ptr=[0,2,3,3]} determines edge ranges
$[0,2)$, $[2,3)$ and $[3,3)$. The edge position indexes the weight;
\code{col\_idx[e]} supplies the source ID. Thus the first row visits source IDs
0 and 1, forms messages 8 and 15 and produces 23. The second produces 6; the
empty third row produces the sum identity 0. Confusing edge positions with
source IDs would already invalidate lowering before any hardware schedule.

\begin{figure}[htbp]
\centering
\begin{tikzpicture}[>=Stealth,every node/.style={font=\small}]
\node[draw,rounded corners,text width=2.2cm,align=center,minimum height=1cm] (domain) {same edge message};
\node[draw,rounded corners,draw=tigateal,thick,text width=4.2cm,align=center,minimum height=1cm,right=1cm of domain] (iter) {select this row's edges};
\draw[->] (domain) -- (iter);
\node[below=0.35cm of iter,text=tigateal] {row 0: edges 0, 1 $\longrightarrow$ sum 23};
\end{tikzpicture}
\caption{Iter introduces coordinate traversal, not a thread assignment.
The highlighted decision is row membership.}
\label{fig:iter-design}
\end{figure}

Dense and generated relations need different coordinate producers. Materializing
all edges merely to resemble CSR can destroy the reason to retain an implicit
relation. Conversely, preserving a constructor in the frontend does not prove
that every backend consumes it implicitly: graph concatenation currently
materializes CSR. Format abstraction and schedule selection must be inspected
on the path actually chosen.

\subsection{Kernel: a concrete, guarded schedule}
The fixture's Kernel operation consumes CSR buffers and field tensors. Its
selected bounded-ragged schedule records these attributes:
\begin{lstlisting}[language=TigaIR]
block_rows = 16, block_neighbors = 16,
num_warps = 4, pipeline_stages = 1
\end{lstlisting}
These are values for this fixture, not global Tiga defaults.
Its execution roles distinguish workgroup destination rows, subgroup
neighbor reduction and scalar lanes. Resource and handoff attributes describe
global reads, private register state, subgroup reduction and global output writes.

The important transition is from enumerating the correct coordinates to deciding
how multiple coordinates share hardware. A maximum-degree guard can admit a
masked bounded-neighbor tile; without that bound, the same launch geometry may
be incorrect or wasteful. A high-degree split requires explicit partial states
and a legal merge. No universal schedule should be inferred from the simple
case's kernel. Provider-neutral scheduling still needs target capabilities and
resource constraints; it does not imply identical machine instructions on every
backend.

\subsection{Task and Storage: dependencies are not residency}
The checked distributed fixture introduces partition and halo planning, then
pack, exchange and unpack events. Interior work depends on partition readiness;
boundary work additionally depends on unpack completion. A join represents the
completion of both branches. This fixture demonstrates the dependency vocabulary;
the public executor uses the communication-then-compute order described in
Section~\ref{sec:runtime}.

Storage represents a logical region and version separately from physical
instances, which carry capacity, memory space, device and layout.
\code{gf\_storage.transfer} produces a readiness event;
\code{gf\_storage.release} depends on completion before storage is recycled;
\code{gf\_storage.join} joins events without specifying a host-wide synchronize.
These operations establish a vocabulary for movement and lifetime, not proof
that a cost model automatically supplies an out-of-core schedule for any kernel.

For a paged graph, correctness requires that the page's topology and field
version agree, transfer completes before consumption, and storage is not reused
until the last consumer finishes. The current public CPU/CUDA paged-CSR executor
uses host-side orchestration for these dependencies. Whole-Tensor
least-recently-used (LRU) eviction is a
separate policy from partitioning one oversized active working set.

\subsection{Transformations, verification and differentiation}
\begin{table}[ht]
\centering\small
\begin{tabularx}{\linewidth}{l Y Y}
\toprule
Transformation & Required fact & Failure if ignored \\
\midrule
Row regrouping & Reducer association/order contract & Different state or numerical behavior \\
Bounded neighbor tile & Degree/shape guards and masked accesses & Truncated rows or out-of-range loads \\
Producer--consumer fusion & Compatible effects and dependencies & Stale values or duplicated side effects \\
Topology reuse & Matching frozen snapshot and bindings & Wrong edges after mutation \\
Storage reuse & Last consumer completed & Use-after-recycle or stale version \\
VJP construction & Saved/recomputed primal state and fixed topology & Incorrect input/parameter gradients \\
\bottomrule
\end{tabularx}
\caption{Legality obligations, not a claim that every possible violation is
statically proven absent. Runtime guards and negative tests remain necessary.}
\end{table}

Differentiation can change the access pattern: source gradients accumulate over
reverse incidence, and edge gradients retain edge identity. A forward schedule
need not be a good backward schedule. The compiler therefore needs a distinction
between semantic VJP generation, schedule selection and runtime saved-state
management. For current Torch CSR replay, that semantic bridge is implemented
through Torch autograd rather than a generated TTIR backward kernel; dedicated
edge-NN and native VJP paths have different contracts.

From the report repository, reproduce the architectural example with an
installed \code{gf-opt} compiler executable:
\begin{lstlisting}[language=TigaShell]
gf-opt generated/ir/distributed-plan.mlir \
  -gf-lower-domain-to-iter \
  -gf-lower-iter-to-kernel -gf-select-kernel-schedule \
  -gf-plan-distributed-tasks
\end{lstlisting}
The \path{generated/ir/} directory contains the individual stages in
\code{domain.mlir}, \code{iter.mlir}, \code{kernel.mlir} and \code{task.mlir}.
The artifact exporter checks them against compiler output before copying them.
